\section{A SOLUÇÃO PROPOSTA: APLICATIVO WeFIND}

A ideia de criar a plataforma WeFIND surgiu a partir de um episódio real vivenciado no círculo pessoal do autor, no qual o desaparecimento de um animal de estimação evidenciou, na prática, as graves limitações de depender apenas de grupos em redes sociais e cartazes de rua. Durante os esforços de busca, constatou-se a sobrecarga enfrentada pela tutora: a obrigação de realizar postagens diárias em múltiplos grupos fragmentados por bairros no Facebook, a necessidade constante de retificar e atualizar dados importantes esquecidos nas publicações iniciais e o retrabalho manual de compartilhamento em suas próprias redes sociais. Como desfecho desse caso, o animal infelizmente não foi localizado, o que demonstrou que a demora na comunicação, a dispersão dos dados e o uso de redes sociais comuns diminuem consideravelmente as chances de resgate nos momentos mais críticos.

A partir dessa constatação empírica, a observação estendeu-se ao comportamento de grupos virtuais em diversas cidades brasileiras, confirmando que a desarticulação das buscas é generalizada. Na dinâmica de redes sociais comuns (como grupos de Facebook, conversas de WhatsApp e postagens no Instagram), os relatos comunitários sofrem com a falta de padronização visual dos pets, com a dificuldade de consolidar novas pistas em uma linha do tempo confiável e com a sobrecarga de quem tenta organizar informações desencontradas entre dezenas de comentários paralelos. O desgaste emocional provocado por essa fragmentação evidenciou a necessidade urgente de uma plataforma exclusivamente desenhada para a causa animal, imune aos ruídos e distrações dos feeds generalistas.

Para solucionar esses entraves, o WeFIND foi concebido como uma aplicação colaborativa integrada, voltada a centralizar o cadastro, a busca georreferenciada e a divulgação de animais em três situações fundamentais: Perdidos, Encontrados e Para Adoção. O projeto apoia-se em quatro pilares conceituais principais:

O primeiro pilar compreende a busca espacial e geolocalização ativa em tempo real. O mapa interativo constitui o núcleo visual do sistema, permitindo que qualquer pessoa visualize animais desaparecidos ou avistados nas imediações de onde transita, configure o raio geográfico de interesse e requisite o traçado dinâmico de rotas viárias com estimativa de tempo de deslocamento até o ponto do registro. Em complemento, o feed contínuo organiza os anúncios ordenados pela distância física até o usuário, priorizando publicações no bairro e na vizinhança imediata.

O segundo pilar estabelece o rastreamento colaborativo e engajamento comunitário ágil. Reconhecendo que animais perdidos circulam continuamente pelas vias públicas, a aplicação institui o mecanismo de avistamento comunitário, no qual qualquer cidadão que veja um pet na rua pode fotografá-lo e demarcar o local exato no mapa, gerando notificações imediatas ao tutor e atualizando o rastro geográfico de busca em uma linha do tempo transparente.

O terceiro pilar fundamenta-se na prevenção e na guarda responsável. Por meio do cadastro e identificação de animais tutelados, o tutor registra preventivamente seus animais sob guarda responsável com fotos, marcas características e histórico veterinário, obtendo um documento de identificação digital provido de código QR exclusivo. Esse código pode ser fixado diretamente na coleira do animal, possibilitando que qualquer pessoa que o encontre escaneie a identificação e acione o responsável imediatamente, sem depender de aparelhos específicos de leitura de microchips em clínicas veterinárias.

Por fim, o quarto pilar protege a privacidade, a integridade e a segurança dos usuários. Em conformidade com as diretrizes da Lei Geral de Proteção de Dados (LGPD), o WeFIND elimina a obrigatoriedade de exibição pública de telefones ou endereços residenciais exatos. Toda a conversa preliminar ocorre em um canal de bate-papo privativo em tempo real dentro da própria aplicação, protegendo tutores contra trotes e golpes oportunistas.

\subsection{Justificativa da Abordagem Móvel em Relação à Web}

A decisão de desenvolver o WeFIND como um aplicativo para celular, em vez de um site comum de internet, fundamentou-se nas exigências práticas de agilidade, mobilidade e prontidão necessárias em situações de emergência animal nas ruas. Um site para computador ou navegador traz a vantagem de não precisar ser instalado, mas tem uma grande limitação: ele depende de a pessoa lembrar-se de abrir o navegador e carregar a página várias vezes ao longo do dia, algo que raramente acontece na rotina das pessoas.

Essa exigência de mobilidade reflete-se na própria dinâmica espacial das ocorrências. Conforme documentado no projeto \textit{FugaPet-Rotas} (SOUZA et al., 2022), animais desorientados em vias públicas possuem grande capacidade de locomoção diária, demandando que a resposta comunitária aconteça exatamente no ambiente físico em que o evento se desenrola: nas calçadas, praças e vias de trânsito, onde cidadãos e voluntários circulam portando smartphones conectados.

Diante dessa dinâmica urbana, o ecossistema móvel destacou-se como a abordagem indispensável por quatro razões práticas:

Primeiramente, pela capacidade de disparo de notificações sonoras e vibratórias instantâneas (\textit{push notifications}). Enquanto uma página web fechada permanece silenciosa e inativa, o aplicativo instalado no smartphone é capaz de alertar o usuário no bolso no exato momento em que um animal foge a poucas quadras de sua posição geográfica atual, mobilizando a comunidade em tempo hábil.

Em segundo lugar, pelo acesso nativo aos sensores de geolocalização por satélite (GPS) e à câmera fotográfica de alta resolução. Em trânsito nas calçadas ou praças, o informante consegue registrar um avistamento em segundos, fotografando o animal e capturando as coordenadas geográficas precisas do ponto do avistamento com um único toque, eliminando a digitação manual de endereços.

Em terceiro lugar, pela viabilidade da guarda preventiva offline. As características físicas, registros de vacinas e fotos do pet permanecem sincronizados localmente no smartphone do tutor. Em eventual fuga acidental, o alerta de desaparecimento é acionado de maneira imediata a partir do cadastro do animal tutelado, evitando que o tutor sob forte abalo emocional perca tempo procurando fotos antigas ou redigitando características.

Por fim, para alcançar também as pessoas que ainda não possuem o aplicativo instalado no aparelho, o WeFIND integra a geração automática de cartazes digitais diagramados com código QR escaneável. Dessa forma, qualquer pessoa na rua que aponte a câmera do celular para o cartaz impresso em postes ou comércios tem acesso imediato à publicação, facilitando o apoio da comunidade sem barreiras técnicas.

Com a concepção, os pilares e a abordagem da plataforma WeFIND fundamentados, o Capítulo 6 apresenta as tecnologias e a modelagem do banco de dados relacional que sustentam o funcionamento da aplicação.
\clearpage
